\section{The finite input and its geometric checker}\label{sec:certificates}

This section specifies the 481-phase input and proves the geometric
upper bounds computed by its checker. The checker evaluates whole
closed root cells and whole angular intervals. Its input is an integer
field table and a deletion-rank table; neither sampled risks nor
search priorities are acceptance premises. The concrete arithmetic
results and execution record are given in Section~\ref{sec:numerics}.

\subsection{A schedule defined by reverse deletion ranks}\label{sec:rankinput}
Index the four types by \(r=0,1,2,3\), with
\(\typ(p)=(p_x\bmod2)+2(p_y\bmod2)\). The file
\texttt{field\_units.txt} contains 1024 nonnegative integers
\(u_{r,a,b}\), in that order, with \(0\le a,b<16\).
Set \(G=256\) and write
\(C_{ij}=(i/G,(i+1)/G)\times(j/G,(j+1)/G)\). Define
\begin{equation}\label{eq:rankfield}
 W_r^{(0)}\big|_{C_{ij}}
 =\frac{249}{250\cdot10^8}
       u_{r,\lfloor i/16\rfloor,\lfloor j/16\rfloor}.
\end{equation}
All fields and masks are zero on every fine-grid boundary and outside
the square. Thus the coarse table defines the values, but the boundary
convention is that of the fine grid.

The file \texttt{full\_rank.bin} contains \(4G^2=262144\)
little-endian signed 32-bit integers \(\nu_{r,i,j}\), in type, first
coordinate, second coordinate order. A zero-valued cell has rank
\(-1\); every positive cell has a rank in \(\{0,\ldots,480\}\).
For \(0\le j\le481\), let
\[
 W_r^{(j)}=W_r^{(0)}\ind_{\{\nu\ge j\}},
 \qquad E_r^{(j)}=W_r^{(0)}\ind_{\{\nu=j\}}
       \quad(0\le j<481).
\]
The indicators here are type-dependent cell indicators.
Then \(W^{(j)}-W^{(j+1)}=E^{(j)}\), and \(W^{(481)}=0\).
For every nonempty rank class set
\(\delta_j=\max_{r,i,k}E_r^{(j)}|_{C_{ik}}>0\).
Use the phases in decreasing rank order, from 480 to 0. On the phase
of duration \(\delta_j\), use the mask
\begin{equation}\label{eq:rankmask}
 f_r=E_r^{(j)}/\delta_j.
\end{equation}
It has rational values in \([0,1]\). This defines the chronological
schedule completely, including its durations and masks.

\begin{lemma}[Reverse deletion and the common deadline]\label{lem:reverse}
The field at the end of the rank-\(j\) phase is \(W^{(j)}\).
Every positive typed cell is used once, and its deadline is the end
of its rank's phase. Consequently, a risk bound for \(W^{(j)}\)
on every rank-\(j\) root cell is a bound for that cell's actual
common-deadline risk. The terminal field is \(W^{(0)}\), and
\[
 A=\frac14\sum_r\int W_r^{(0)},\qquad
 T=\sum_{j=0}^{480}\delta_j.
\]
\end{lemma}
\begin{proof}
At the indicated endpoint the phases already run are exactly those
with ranks \(k\ge j\). Their exposures sum to
\(\sum_{k=j}^{480}E^{(k)}=W^{(j)}\). This equality holds for all
four fields at the same time. A positive typed cell belongs to just
one rank class, so its last use is precisely that phase. Finally the
classes partition all positive typed cells. There are no overlapping
accepted-region areas to count with multiplicity.
\end{proof}

All intermediate fields are componentwise at most \(W^{(0)}\).
Capacities and squared-chord risk are monotone for nonnegative fields.
In particular, a terminal-field capacity bound controls the entire
schedule. The rank correspondence itself is checked before any
geometric calculation: the program rejects an out-of-range rank, a
mismatch between positivity and rank, or an empty requested phase.

\subsection{Exact maxima on a closed root cell}
Let \(F\) be a nonnegative rectangular step field, extended by zero
outside the unit square. For a direction \((1,u)\), \(u>0\), put
\[
 C(v;u)=\int_{\R}F(t,v+ut)\dd t.
\]
For a closed root box \(B=[a_0,a_1]\times[b_0,b_1]\), its possible
intercepts form the interval
\(l(u)=b_0-ua_1\le v\le h(u)=b_1-ua_0\).
At each field vertex \((a,b)\), combine the signed corner
coefficients of its incident rectangles, with exterior values zero.
The coefficients for one rectangle of value \(f\), in lower-left,
lower-right, upper-left, upper-right order, are \(-f,+f,+f,-f\).

\begin{lemma}[Corner representation]\label{lem:corners}
For \(u>0\),
\begin{equation}\label{eq:corners}
 C(v;u)=\frac1u\sum_{a,b}k_{a,b}(v-b+ua)_+.
\end{equation}
\end{lemma}
\begin{proof}
For the indicator of \((a_0,a_1)\times(b_0,b_1)\), integrate
\(\ind_{v+ut>b_0}-\ind_{v+ut>b_1}\) over \(a_0<t<a_1\).
For the first indicator the result is
\([(v+ua_1-b_0)_+-(v+ua_0-b_0)_+]/u\), and similarly for the
second. Add the identities with the rectangle weights and combine
shared corners. A nonaxis line meets the zeroed grid lines in finitely
many points, which do not affect this integral.
\end{proof}

\begin{proposition}[Exact closed-box envelope]\label{prop:envelope}
For fixed \(u>0\), the maximum chord of all roots in \(B\) is
\begin{equation}\label{eq:envelope}
 U_B(u)=\max\bigl\{C(v;u):
 v\in\{l(u),h(u)\}\cup
       \{b-ua:l(u)\le b-ua\le h(u)\}\bigr\}.
\end{equation}
The second set ranges over the field vertices. In particular this
maximum is obtainable by exact rational arithmetic at rational \(u\).
\end{proposition}
\begin{proof}
By \eqref{eq:corners}, the chord is continuous and piecewise affine
in \(v\), with knots among \(b-ua\). The maximum on an interval is
at an endpoint or at an included knot. Every possible root intercept
belongs to that interval, and every intercept in it is attained by
some root of the box.
\end{proof}

The implementation sorts projected vertices, combines equal knots,
and accumulates the slope and then the value of the chord numerator.
A segment tree returns the largest included knot value; the two
interval endpoints are evaluated separately. The zero value and zero
slope beyond the last knot are checked. Taking a maximum over all
knots, rather than only the root interval, also gives a valid bound
simultaneously for every root in the square.

\subsection{A rigorous angular-bin bound by vertical dilation}
For \(\delta\ge0\), let
\[
 (D_\delta F)(t,y)=\sup_{|z-y|\le\delta}F(t,z).
\]
\begin{lemma}[Tube domination]\label{lem:tube}
For \((x_0,y_0)\in[0,1]^2\) and \(|u-u_0|\le\delta\),
\[
 c_{F,1}((x_0,y_0),u)
 \le c_{D_\delta F,1}((x_0,y_0),u_0).
\]
\end{lemma}
\begin{proof}
Only \(0\le t\le1\) contributes. On this interval,
\(|(u-u_0)(t-x_0)|\le\delta\). The point of the original line is
therefore among those in the vertical supremum at the corresponding
point of the central line. Integrate the resulting pointwise bound.
\end{proof}

Choose \(N=2048\), \(H=2N\), \(\delta=1/H\), and
\(u_k=(2k+1)/H\), \(0\le k<N\). Lemma~\ref{lem:tube} gives
\begin{equation}\label{eq:tube}
 \sup_{x\in B}\int_0^1 c_{F,1}(x,u)^2\dd u
 \le\frac1N\sum_{k=0}^{N-1}
       \left(\sup_{x\in B}c_{D_\delta F,1}(x,u_k)\right)^2.
\end{equation}
This is a bound on all angles in each bin, not a quadrature
approximation whose error has been omitted.

The dilation is itself a finite rational rectangular field. In each
column, let the original cell values be \(f_j\), with exterior values
zero. Since \(\delta<1/(2G)\), the central strip of cell \(j\)
has value \(f_j\), and the strip of width \(2\delta\) about its
boundary \(j/G\) has value \(\max(f_{j-1},f_j)\). These strips
give an explicit decomposition, including the two outer boundaries.
Values on strip endpoints do not affect the positive-slope chord.
The program forms the integer type-pair field before this dilation;
it does not round each type's field separately.

All dilated vertices have coordinates \((a/H,b/H)\) with integers
\(a,b\). At \(u=p/H\), their projected intercepts are \(Hb-pa\).
The integer sweep computes a numerator whose raw-field chord is
that numerator divided by \(pH\). With rounding scale
\(Q_c=10^6\), the true scaled chord in \eqref{eq:rankfield} is
bounded in \(1/Q_c\) units by
\[
 L=\left\lceil\frac{249\,\text{numerator}}{25000\,pH}\right\rceil.
\]
This uses \(Q_c\,249/(250\cdot10^8)=249/25000\).
The bin contributes the integer upper bound \(\lceil L^2/N\rceil\)
in units \(1/Q_c^2\). All divisions round upward.

\subsection{Complete Farey partitions and rational integral bounds}
If the tube estimate is too large to accept a phase, the checker
uses the same residual field and every root box of that phase in a
second calculation. Its partition is the list of all reduced fractions
in \([0,1]\) with denominator at most \(G=256\).

A projected-knot ordering or a knot's membership in the root intercept
interval can change only at a slope joining two field or root-box
vertices. Those vertices are on the \(1/G\) grid, so every such slope
in \([0,1]\) has reduced denominator at most \(G\). On an open
interval between consecutive fractions, the candidate set in
\eqref{eq:envelope} is fixed and every candidate has the form
\(a+b/u\). Thus \(U_B(u)\) is convex as a function of \(1/u\).
Its endpoint secant in that coordinate bounds it from above.
Switches between maximizing candidates require no extra partition.
At a positive partition endpoint, continuity of the chord and of the
root intercept interval ensures that each one-sided candidate limit
is bounded by the endpoint maximum used in the calculation.

Let the consecutive positive fractions be \(\alpha=a/b\) and
\(\beta=c/d\), so \(bc-ad=1\) and \(\beta-\alpha=1/(bd)\).
Let \(A,B\) be upward-rounded endpoint chord values in units
\(1/Q_c\). The reciprocal-coordinate secant through these rounded
values remains an upper bound for \(Q_c U_B\).

\begin{lemma}[Rational affine integral bound]\label{lem:affine}
An upper bound for \(Q_c^2\int_\alpha^\beta U_B(u)^2\dd u\) is
the smaller of \(\max(A^2,B^2)/(bd)\) and
\begin{equation}\label{eq:affine}
 \begin{cases}
 (A^2+AB+B^2)/(3bd),&B\le A,\\[2pt]
 (A'^2+A'B'+B'^2)/(3bd),&B>A,
 \end{cases}
\end{equation}
where
\[
 A'=A+\frac{(B-A)bc}{(ad+bc)^2},\qquad
 B'=B+\frac{(B-A)ad}{(ad+bc)^2}.
\]
\end{lemma}
\begin{proof}
Write the reciprocal-coordinate secant as \(V(u)=s+t/u\).
It is nonnegative and monotone on the interval. If \(B\le A\),
then \(t\ge0\), so the endpoint secant in \(u\) is an upper bound.
If \(B>A\), then \(t<0\) and \(V\) is concave; its tangent at
\((\alpha+\beta)/2\) is an upper bound. Substitution and
\(bc-ad=1\) give tangent endpoint values \(A',B'\).
The integral of the square of an affine function with endpoint
values \(x,y\) on an interval of length \(1/(bd)\) equals
\((x^2+xy+y^2)/(3bd)\). The constant endpoint-maximum bound is
valid by monotonicity. Take the smaller of these two valid bounds.
\end{proof}

On the first interval \((0,1/G)\), every active candidate is bounded
and of the form \(a+b/u\), hence \(b=0\). Its maximum is constant,
so the first positive endpoint suffices. The value at angle zero
does not affect the angular integral; it does not replace the axis
capacity check below.

For an ordinary vertex \((i/G,j/G)\) at angle \(p/q\), the
projection is \(qj-pi\). The integer numerator is normalized by
\(pG\), and the rounded scaled chord is
\(\lceil249\,\text{numerator}/(25000\,pG)\rceil\).
The tangent values in \eqref{eq:affine} are retained as exact rational
numbers. Only the final integral bound is rounded upward to an integer
in units \(1/Q_c^2\). No logarithmic numerical integration is used.

\subsection{Four direction families, integer ranges and acceptance}
For each of the six unordered type pairs, the checker transforms the
field and root box by each of
\[
 (x,y),\quad(x,1-y),\quad(y,x),\quad(y,1-x).
\]
These give the four direction families of \eqref{eq:risk} and preserve
the major-coordinate integration parameter. Type labels are retained.
No transpose symmetry of the fields or symmetry of the random process
is assumed. Each type receives three pairs and four families.

If their twelve integral bounds sum to the integer \(J\), then
\begin{equation}\label{eq:certRisk}
 \mathcal R_r(W^{(j)};x)\le\frac{19J}{750Q_c^2}
 \quad\hbox{throughout the checked root box}.
\end{equation}
The acceptance test is
\[
 J\le\left\lfloor\frac{9999\cdot750Q_c^2}{190000}\right\rfloor
      =39469736842105.
\]
It implies \(\mathcal R_r\le9999/10000<1\). Each output row is
identified by type and fine-cell coordinates. The replay driver
compares the complete set of rows, and its cardinality, with the rank
class to exclude missing or repeated cells. It independently compares
the row integers to the rational risk threshold. Passing only selected
cells or retaining an unsuccessful tube bound is insufficient.

For this fixed input, the largest single-type integer is 628311458
and the largest type-pair integer is 1181489833. Deletion and dilation
cannot increase the latter bound. Its physical chord, in \(1/Q_c\)
units and rounded upward, is less than 13000000. A corner coefficient
has absolute value at most four times the pair maximum, below
\(2^{33}\). Even at the largest supported angular parameter
\(N=65536\), projected coordinates are below \(2^{35}\); a coarse
sum-of-absolute-values bound puts sweep intermediates below \(2^{86}\).
Actual chord numerators are below \(2^{65}\) before multiplication
by 249. The values and arithmetic used for chord conversion and
squaring are therefore within signed 128-bit range.
The sum over all bins of one tube integral is at most
\(13000000^2+N\); twelve such sums, and their risk numerator after
multiplication by 19, are also far below that range.
The same bound holds for a Farey integral with the number of intervals
in place of \(N\), using the endpoint-maximum bound on each interval.

For the Farey tangent put \(D=(ad+bc)^2\). Since
\(ad,bc\le G^2\), its rational numerator uses
\[
 A_L=AD+(B-A)bc,\qquad B_L=BD+(B-A)ad,
\]
and denominator \(3bdD^2\). The bound
\(A_L,B_L\le13000000(D+G^2)\) puts
\(A_L^2+A_LB_L+B_L^2\) below \(2^{117}\), and the denominator
below \(2^{86}\). The final ceiling addition also fits. Ordinary
integer division and signed 128-bit intermediates suffice for these
fixed data. These estimates are part of the soundness argument;
runtime undefined-behavior checks provide an additional implementation
check. They are not a claim that arbitrary int64 field tables would
be safe without a new range bound.

\subsection{Capacity: axes, finite directions and the infinite tail}
Exchange the two coordinates, transforming the field and root as
well, when the second coordinate of the direction is larger in
absolute value. Reverse the direction's orientation if necessary.
It then has the form \(v=(h,b)\), with \(|b|\le h\). Put
\(q=b/h\); the horizontal-coordinate chord below, divided by
\(2h\), is exactly the capacity in \eqref{eq:capacity}.
For a nonaxis slope \(q\ne0\), the chord of a rectangular cell
\((A,B)\times(C,D)\) at intercept \(v\) is
\[
 \pos{\min\left(B,\max\left(\frac{C-v}{q},\frac{D-v}{q}\right)\right)
 -\max\left(A,\min\left(\frac{C-v}{q},\frac{D-v}{q}\right)\right)}.
\]
It is continuous and piecewise affine in \(v\), with knots among
\(C-qA,C-qB,D-qA,D-qB\). Hence the maximum of the summed chord
is attained at an intercept through a grid vertex. The independent
capacity calculation uses this direct rectangle-intersection formula,
not the corner-sweep routine.

Axis directions need a different rule. In an open horizontal strip
the chord is its weighted row sum, whereas on a zeroed horizontal
boundary it is zero. Check all open-strip values, and use column sums
for the vertical direction. Further boundary zeroing on the fine grid
cannot increase these capacities: it removes finitely many points
from a nonaxis chord and may zero an entire axis chord.

Enumerate all 72 primitive unoriented directions of heights one
through seven. A primitive direction has nonzero parity, so it pairs
the four types into two unordered pairs; there are 144 finite
direction/type-pair checks. The 16 directions of height at most three
are designated. For every remaining direction with \(h\ge8\),
\[
 L_{v,r}(W^{(0)};x)
 \le\frac{\|W_r^{(0)}+W_{r+\bar v}^{(0)}\|_\infty}{2h}
 \le\frac{\|W_r^{(0)}+W_{r+\bar v}^{(0)}\|_\infty}{16}.
\]
Checking the six pair maxima therefore completes the infinite tail.
All these calculations may use the coarse \(16\times16\) envelope.

\begin{proposition}[Finite geometric checking contract]\label{prop:contract}
Suppose the specified field and rank tables satisfy the stated format
and positivity correspondence, every rank class is covered by its
checked closed cells, and the exact tube or Farey calculation proves
the common bound \(U_R\le9999/10000\). Suppose the capacity
calculation covers all finite directions and the tail, with bounds
\(U_H<K\) and \(U_L<\lambda<1-\rho\). If \(9999/10000<d<1\),
then the schedule \eqref{eq:rankmask} has the geometric input
properties \eqref{eq:heavyinput}--\eqref{eq:riskinput}, with uniform
positive margins and the exact mass given by its finite field sum.
\end{proposition}
\begin{proof}
The capacity calculation covers every line intercept and direction,
and terminal-field domination transfers it to all times.
Proposition~\ref{prop:envelope}, Lemma~\ref{lem:tube} and
Lemma~\ref{lem:affine} justify the angular bounds on every checked
closed cell. Equation~\eqref{eq:certRisk} gives the risk bound, while
Lemma~\ref{lem:reverse} identifies the four tested fields with the
four actual fields at that root's deadline. The finite maxima leave
the positive margins \(K-U_H\), \(\lambda-U_L\) and
\(d-9999/10000\), independent of \(n\). The rank classes partition
all positive typed cells, giving the mass without overlap ambiguity.
Proposition~\ref{prop:pair} and \eqref{eq:discrete481} then provide
the uniform discrete comparison used by the stochastic argument.
\end{proof}

The executable \texttt{independent/check\_phase.cpp}, with
\texttt{farey\_independent.inc}, implements these two angular bounds.
Its explicit rectangles, standard sorting, segment tree, ordinary
division, wider intermediates and four separate coordinate transforms
provide a second implementation of the geometric conditions.
Both implementations still use the mathematical lemmas proved here.
Input hashes identify the data; they do not prove the inequalities.
For a complete reproduction the checker is rebuilt and every phase
is evaluated from the fixed input, followed by per-cell coverage and
integer comparisons. Reusing cached phase reports without binding them
to the current inputs and checker is not such a reproduction.
